\begin{figure}[tb]
\centering
\begin{tikzpicture}
\begin{axis}[
 width=.82\linewidth,height=.54\linewidth,
 xmin=0,xmax=6,ymin=0,ymax=5,
 xlabel={normalized torque-per-loss coordinate},
 ylabel={field-weakening coordinate},ticks=none,axis lines=left,clip=false]
 \addplot[name path=low,estimated quantity,strong line,domain=0:5.7] {1.15+4.5/(x+1.25)};
 \addplot[name path=base,reference quantity,strong line,domain=.4:5.8] {.58*x+.35};
 \addplot[name path=high,highlighted quantity,strong line,domain=.3:5.8] {4.55-.33*x};
 \addplot[name path=floor,draw=none,domain=0:6] {0};
 \addplot[derived muted quantity,derived muted region]
   coordinates {(2.55,2.91) (3.25,2.24) (4.35,2.89) (3.67,3.34)} \closedcycle;
 \node[estimated quantity,anchor=west,compact text] at (axis cs:4.1,2.05) {low-speed torque};
 \node[reference quantity,anchor=west,compact text,rotate=26] at (axis cs:3.7,2.7) {base speed};
 \node[highlighted quantity,anchor=west,compact text,rotate=-15] at (axis cs:.7,4.45) {maximum speed};
 \node[derived muted quantity,align=center,annotation text] at (axis cs:3.45,2.75)
  {feasible\\design window};
 \addplot[only marks,mark=*,neutral quantity] coordinates {(3.35,2.7)};
 \draw[coordinate axis,standard line] (axis cs:3.35,2.7)--(axis cs:2.95,2.2)
  node[below,compact text] {active margin};
\end{axis}
\end{tikzpicture}
\caption{Multiple operating points as intersecting design regions in a
two-dimensional normalized slice.  Each boundary is a level set of the same
current-allocation value function.  Only the common derived quantity window satisfies the
entire torque-speed specification.}
\label{fig:multipoint-window}
\end{figure}
